\documentclass[aspectratio=169, xcolor={table,xcdraw}]{beamer}

\usetheme{Madrid}
\usecolortheme{default}
\definecolor{ucbblue}{RGB}{0, 51, 102}

\setbeamercolor{structure}{fg=ucbblue}
\setbeamercolor*{palette primary}{fg=white,bg=ucbblue}
\setbeamercolor*{palette secondary}{fg=white,bg=ucbblue!80}
\setbeamercolor*{palette tertiary}{fg=white,bg=ucbblue!90}
\setbeamercolor{title}{fg=white, bg=ucbblue}
\setbeamercolor{frametitle}{fg=white, bg=ucbblue}

\usepackage[T1]{fontenc}
\usepackage[utf8]{inputenc}
\usepackage{tgpagella}
\usefonttheme{serif}
\usepackage[spanish, es-noshorthands]{babel}
\usepackage{amsmath, amssymb}
\usepackage{booktabs}
\usepackage{tabularx}
\usepackage[most]{tcolorbox}
\usepackage[american]{circuitikz}

\title[Op-Amps: Ideal a Real]{Del Modelo Ideal al Circuito Físico: Amplificador Diferencial y Errores DC}
\subtitle{Circuitos Electrónicos II (IMT-221) - Semana 7, Sesión 2}
\author[B. Quiroga Turdera]{M.Sc. Ing. Bernardo Quiroga Turdera}
\institute[UCB Tarija]{Universidad Católica Boliviana ``San Pablo'' — Sede Tarija}
\date{}

\begin{document}

\begin{frame}
    \titlepage
\end{frame}

\begin{frame}{Pregunta Experimental Fundamental}
    \begin{tcolorbox}[colback=ucbblue!5, colframe=ucbblue, title=\textbf{El Problema de Ingeniería[cite: 11]}]
        ¿Por qué un amplificador diferencial físico puede entregar una salida distinta de la predicción ideal y cómo puede diagnosticarse y reducirse esa diferencia?[cite: 11]
    \end{tcolorbox}
    
    \vspace{0.4cm}
    \textbf{Secuencia de Trabajo (Hacia Lab 3)[cite: 11]:}
    \begin{center}
        Predicción Ideal $\rightarrow$ Modelo Real Mínimo $\rightarrow$ Predicción de Error $\rightarrow$ Simulación $\rightarrow$ \\
        Plan de Medición $\rightarrow$ \textit{Implementación Física (Sem 8)} $\rightarrow$ \\
        Comparación $\rightarrow$ Diagnóstico $\rightarrow$ Compensación[cite: 11]
    \end{center}
\end{frame}

\begin{frame}{Recuperación: Amplificador Diferencial Ideal}
    \begin{columns}
        \begin{column}{0.4\textwidth}
            \begin{circuitikz}[scale=0.8, transform shape, thick]
                \draw (0,0) node[op amp] (opamp) {};
                \draw (opamp.+) to[R, l_=$R_3$] ++(-2,0) node[left]{$V_2$};
                \draw (opamp.+) to[R, l=$R_4$] ++(0,-2) node[ground]{};
                \draw (opamp.-) to[R, l_=$R_1$] ++(-2,0) node[left]{$V_1$};
                \draw (opamp.-) -- ++(0,1.5) coordinate(tmp) to[R, l=$R_2$] (tmp -| opamp.out) -- (opamp.out) to[short, -o] ++(0.5,0) node[right]{$V_o$};
            \end{circuitikz}
        \end{column}
        \begin{column}{0.6\textwidth}
            \small
            \textbf{Análisis Nodal (Modelo Ideal)[cite: 11]:} \\
            Por divisor de tensión en la entrada no inversora ($i_+ = 0$):
            \begin{equation*}
                V_+ = V_2\frac{R_4}{R_3+R_4} \text{[cite: 11]}
            \end{equation*}
            Por realimentación negativa: $V_x = V_- \approx V_+$[cite: 11]. \\
            KCL en nodo inversor ($i_- = 0$):
            \begin{align*}
                \frac{V_1 - V_x}{R_1} + \frac{V_o - V_x}{R_2} &= 0 \text{[cite: 11]} \\
                \frac{V_o}{R_2} &= V_x\left(\frac{1}{R_1} + \frac{1}{R_2}\right) - \frac{V_1}{R_1} \text{[cite: 11]} \\
                V_o &= V_x\left(1+\frac{R_2}{R_1}\right) - V_1\frac{R_2}{R_1} \text{[cite: 11]}
            \end{align*}
        \end{column}
    \end{columns}
\end{frame}

\begin{frame}{Condición de Diferencia Ideal}
    Sustituyendo $V_x = V_2\frac{R_4}{R_3+R_4}$ en la ecuación anterior[cite: 11]:
    \begin{equation*}
        \boxed{ V_o = \left(1+\frac{R_2}{R_1}\right)\frac{R_4}{R_3+R_4}V_2 - \frac{R_2}{R_1}V_1 } \text{[cite: 11]}
    \end{equation*}
    
    Para obtener una resta diferencial pura, forzamos la igualdad de razones resistivas[cite: 11]:
    \begin{equation*}
        \text{Si } \frac{R_2}{R_1} = \frac{R_4}{R_3} = k \implies \boxed{ V_o = k(V_2 - V_1) } \text{[cite: 11]}
    \end{equation*}
    
    \textbf{Ejemplo Pedagógico (Diseño Base Lab 3)[cite: 11]:} $R_1=R_3=10\,\text{k}\Omega$, $R_2=R_4=100\,\text{k}\Omega \implies k=10$[cite: 11]. \\
    Si $V_1=0.20\,\text{V}$ y $V_2=0.25\,\text{V}$, $V_{o,\mathrm{ideal}} = 10(0.25 - 0.20) = \boxed{0.50\,\text{V}}$[cite: 11].
\end{frame}

\begin{frame}{Del Modelo Ideal al Op-Amp Real}
    El modelo ideal elimina imperfecciones. Clasificamos los errores físicos[cite: 11]:
    
    \vspace{0.3cm}
    \begin{table}[]
        \centering
        \small
        \renewcommand{\arraystretch}{1.3}
        \begin{tabularx}{\textwidth}{>{\raggedright\arraybackslash\hsize=1.2\hsize}X >{\raggedright\arraybackslash\hsize=0.8\hsize}X >{\raggedright\arraybackslash\hsize=1.0\hsize}X}
        \toprule
        \rowcolor[HTML]{EFEFEF}
        \textbf{Modelo Ideal} & \textbf{Situación Real (Lab 3)} & \textbf{Consecuencia Observable} \\ \midrule
        Diferencia de entrada idealmente nula ($V_+ = V_-$) & $V_{OS}$ (Input Offset Voltage)[cite: 11] & Offset de salida[cite: 11] \\
        $i_+ = i_- = 0$ & $I_{B+}, I_{B-}$ (Bias Currents)[cite: 11] & Caídas adicionales en resistencias[cite: 11] \\
        Corrientes de entrada idénticas & $I_{OS} \neq 0$ (Offset Current)[cite: 11] & Error residual con red balanceada[cite: 11] \\
        Razones resistivas exactas & Tolerancia / Desajuste[cite: 11] & Fuga de modo común ($CMRR$ pobre)[cite: 11] \\
        \bottomrule
        \end{tabularx}
    \end{table}
\end{frame}

\begin{frame}{1. Tensión de Offset de Entrada ($V_{OS}$)}
    \begin{columns}
        \begin{column}{0.4\textwidth}
            \begin{circuitikz}[scale=0.8, transform shape, thick]
                \draw (0,0) node[op amp] (opamp) {};
                \draw (opamp.+) to[V, v=$V_{OS}$, invert] ++(-1.5,0) to[short, -o] ++(-0.5,0) node[left]{$\text{Entrada Real}$};
                \draw (opamp.-) to[short, -o] ++(-2,0);
                \node at (0, -2.0) {\textit{Modelo Equivalente[cite: 11]}};
            \end{circuitikz}
        \end{column}
        \begin{column}{0.6\textwidth}
            \raggedright
            $V_{OS}$ es una pequeña tensión DC equivalente en la entrada necesaria para anular la salida cuando las entradas físicas son cero[cite: 11]. \\
            
            \vspace{0.2cm}
            \textbf{Propagación a la Salida[cite: 11]:} \\
            Este error se multiplica por la "Ganancia de Ruido" ($A_N$)[cite: 11].
            \begin{equation*}
                V_{o,\mathrm{offset}} \approx A_N V_{OS} \text{[cite: 11]}
            \end{equation*}
            \textit{Ejemplo Pedagógico:} Si $A_N = 11$ y $V_{OS} = 1\,\text{mV} \implies V_{o,\mathrm{offset}} \approx 11(1\,\text{mV}) = \boxed{11\,\text{mV}}$[cite: 11].
        \end{column}
    \end{columns}
\end{frame}

\begin{frame}{2. Corrientes de Polarización ($I_B$) y Offset ($I_{OS}$)}
    Un Op-Amp real requiere corrientes DC minúsculas para polarizar sus transistores de entrada[cite: 11].
    \begin{equation*}
        I_B = \frac{|I_{B+}|+|I_{B-}|}{2}, \quad I_{OS} = |I_{B+} - I_{B-}| \text{[cite: 11]}
    \end{equation*}
    
    \textbf{Mecanismo de Error[cite: 11]:} \\
    Estas corrientes atraviesan las resistencias equivalentes de la red, generando caídas de tensión parásitas ($V_{err} = I_B R_{eq}$)[cite: 11]. \\
    
    \vspace{0.2cm}
    \textit{Ejemplo Pedagógico:} Si $I_B = 100\,\text{nA}$ y $R_{eq} = 100\,\text{k}\Omega \implies V_{err} = (100\,\text{nA})(100\,\text{k}\Omega) = \boxed{10\,\text{mV}}$[cite: 11].
    
    \vspace{0.2cm}
    \textbf{Mitigación:} Igualar las resistencias equivalentes vistas por ambas entradas minimiza el efecto de $I_B$, pero el error debido a $I_{OS}$ persistirá[cite: 11].
\end{frame}

\begin{frame}{3. Desajuste de Resistores (Mismatch)}
    La resta diferencial depende de que $\frac{R_2}{R_1} = \frac{R_4}{R_3}$. Las tolerancias rompen esta igualdad[cite: 11].
    
    \textbf{Ejemplo Analítico[cite: 11]:}
    $R_1=R_3=10\,\text{k}\Omega$, $R_2=100\,\text{k}\Omega$, \textbf{$R_4=101\,\text{k}\Omega$}[cite: 11]. \\
    Apliquemos una entrada de \textbf{modo común}: $V_1=V_2=1\,\text{V}$[cite: 11].
    
    Contribución de $V_2$[cite: 11]: 
    \begin{align*}
        V_o(V_2) &= \left(1+\frac{100}{10}\right)\frac{101}{10+101}(1\,\text{V}) \text{[cite: 11]} \\
                 &= 11\left(\frac{101}{111}\right) \approx 10.009\,\text{V} \text{[cite: 11]}
    \end{align*}
    Contribución de $V_1$: $V_o(V_1) = -\frac{100}{10}(1\,\text{V}) = -10\,\text{V}$[cite: 11]. \\
    \textbf{Resultado:} $V_o \approx 10.009 - 10 = \boxed{9\,\text{mV}}$ (Idealmente debería ser $0\,\text{V}$)[cite: 11].
\end{frame}

\begin{frame}{Lógica de Compensación y LTSpice}
    \textbf{Simulación (LTSpice)[cite: 11]:} \\
    Constituye predicción basada en modelo. $e_{\mathrm{sim}} = V_{o,\mathrm{sim}} - V_{o,\mathrm{ideal}}$[cite: 11]. Un error nulo en simulación solo indica coherencia teórica, no prueba el circuito físico[cite: 11].
    
    \vspace{0.4cm}
    \textbf{Flujo de Diagnóstico y Compensación en Laboratorio[cite: 11]:}
    \begin{center}
        $\boxed{\text{Baseline}} \rightarrow \boxed{\text{Cuantificar Error}} \rightarrow \boxed{\text{Descartar Fallas (Wiring)}} \rightarrow$ \\ $\boxed{\text{Compensar}} \rightarrow \boxed{\text{Re-Medir}} \rightarrow \boxed{\text{Comparar}}$[cite: 11]
    \end{center}
    
    Métrica de mejora: $\Delta |e| = |e_{before}| - |e_{after}| > 0$[cite: 11]. \\
    \textit{Nunca compense antes de verificar: alimentación, tierra, pinout, y lazos de realimentación.[cite: 11]}
\end{frame}

\begin{frame}{Puente a Lab 4: Limitaciones Dinámicas}
    Una vez controlados los errores DC (Lab 3), el circuito enfrentará restricciones dinámicas[cite: 11].
    
    \begin{itemize}
        \item \textbf{Ancho de Banda Finito (GBW):} El Op-Amp real no mantiene su ganancia ideal de lazo abierto a cualquier frecuencia[cite: 11].
        \item \textbf{Slew Rate (SR):} La salida física no puede variar arbitrariamente rápido[cite: 11].
        \begin{equation*}
            SR = \max\left|\frac{dV_o}{dt}\right| \text{[cite: 11]}
        \end{equation*}
    \end{itemize}
    \textit{Pregunta Puente:} ¿Por qué un amplificador que funciona en DC produce salida distorsionada al subir la frecuencia o la amplitud?[cite: 11] \textbf{$\rightarrow$ Hacia Lab 4.[cite: 11]}
\end{frame}

\end{document}
